%---------------------------------------------------------------------
\chapter{Introduction}
%---------------------------------------------------------------------

\section{Background}
India regulates electronic waste through the E-Waste (Management) Rules, 2022, notified by the
Ministry of Environment, Forest and Climate Change (MoEFCC) as G.S.R.~801(E) on 2~November
2022~\cite{rules2022}. The rules place an Extended Producer Responsibility (EPR) obligation on
every producer of covered electrical and electronic equipment (EEE): a business that sells
equipment under its brand, or imports it, must ensure that a stated share of the equipment it
placed on the market earlier is recycled. Schedule~III sets the share at 60\% for 2023--24 and
2024--25, 70\% for 2025--26 and 2026--27, and 80\% from 2027--28, computed on a historical sales
cohort determined by product life rather than on current sales~\cite{rules2022}. Producers
register on the Central Pollution Control Board (CPCB) EPR portal and meet the obligation with EPR
certificates that registered recyclers generate for recovered material and transfer to them;
Rules~14 and~15 describe how certificates are adjusted against liability, with earlier liability
taking priority~\cite{rules2022}. The rules were amended in 2023 and 2024; the March 2024
amendment (G.S.R.~164(E)) allows certain return timelines to be relaxed by
order~\cite{amend2024}.

For a producer, this creates a recurring reconciliation problem. The obligation, the sales
history behind it, the certificates bought, the payments made, the transfers recorded on the
portal and the balances the portal shows are different facts held in different documents. CPCB's
own guidance separates them: a commercial agreement or a payment does not prove that a
certificate was transferred, and certificate metal values, equipment mass and money are separate
quantities~\cite{framework,faq2024}. A compliance officer preparing a return must therefore
trace each claimed reduction of the obligation back to evidence, decide whether the evidence
supports it, and record why.

\section{Problem Statement}
A compliance officer at a producer has a dossier for one reporting period: registration
records, obligation statements, sales and product mappings, purchase and payment records,
certificate records and observations, transfer events, portal balances, and the applicable
rules. Some of these records may be inconsistent, duplicated, expired, attributed to the wrong
entity, expressed in the wrong unit or period, unreadable, or may even contain text that tries to
instruct an automated reader. The officer needs to know which part of the obligation is supported,
which discrepancies exist, and what remains uncertain, with every statement traceable to a page
or record. \textbf{The decision supported is: for this producer and period, which obligation
balances are supported by verified evidence, which discrepancies must be followed up, and what
must a human reviewer confirm before the working paper is used?}

\section{Motivation and Why Agentic AI}
Two simpler designs were considered. A purely rule-based reconciler can compute balances once the
records are clean and linked, and ReCycleGuard does contain one, because arithmetic over
compliance quantities must be exact and reproducible. But the investigation that precedes the
arithmetic is open-ended: which record to inspect next depends on what the previous inspection
showed, a discrepancy may need a rule lookup for the period in question, and a weak or ambiguous
result must lead either to another check or to an explicit abstention. A single prompt to a
language model, on the other hand, can neither inspect a dossier of thousands of observations
within its context nor be trusted to compute balances or cite pages it has not read.

The problem therefore suggests an agent in a specific, bounded sense: a loop that (i) chooses the
next check from a fixed set of audited tools, (ii) observes structured evidence, (iii) assesses
whether the evidence supports, refutes or does not settle a candidate finding, and (iv) stops with a
recorded reason, under hard budgets. Retrieval is needed because the rules and the dossier
text are long; human review is needed because the output has regulatory weight. The project also
treats the agentic hypothesis as something to be tested rather than assumed: model-free controls run
through the same pipeline, so the model's own contribution can be measured.

\section{Objectives}
\begin{enumerate}[itemsep=2pt]
\item To design and implement an evidence pipeline that stores originals immutably, extracts
observations with page and row locators, and reconciles obligations with exact decimal
arithmetic.
\item To implement a bounded LangGraph investigation agent with audited tools, a deterministic
finding gate and a human review step bound to an exact report revision.
\item To build a 48-dossier synthetic benchmark with sixteen fault families and an evaluator whose
answers never enter the runtime, prompts or user interface.
\item To measure, against model-free controls and with confidence intervals, how much a locally
served 9-billion-parameter model contributes to correct findings.
\item To test, under a pre-registered protocol, whether a multi-agent design with a retrieval
agent and a classification agent recovers judgement faults that the single-agent loop misses,
and to run the held-out dossiers once.
\end{enumerate}

\section{Scope of the Project}
Included: one producer and one reporting period per dossier; structured records and PDFs;
a curated corpus of two public CPCB documents (the 2022 Rules and the January 2024 FAQ); a local
web interface and HTTP interface for review; local inference only. Excluded: filing on the CPCB
portal, trading or transferring certificates, any statement of legal compliance or fraud, and
real producer data. The output is an advisory working paper that a person must review; the system
never approves its own findings.

\section{Organisation of the Report}
Chapter~2 reviews how EPR reconciliation is done today and the agentic techniques used here.
Chapter~3 states the requirements and data sources, and Chapter~4 presents the architecture and
the multi-agent design. Chapter~5 maps the system to the course's agentic concepts with evidence,
and Chapter~6 describes the implementation and the evaluation harness. Chapter~7 covers tracing,
the benchmark, the metrics and the results, and Chapter~8 the pre-registered multi-agent
experiment. Chapter~9 discusses limitations, failure modes and responsible use, and Chapter~10
concludes.
